\uuid{kzpo}
\exo7id{4715}
\titre{exo7 4715}
\auteur{quercia}
\organisation{exo7}
\datecreate{2010-03-16}
\isIndication{false}
\isCorrection{true}
\chapitre{Suite}
\sousChapitre{Suite définie par une relation de récurrence}
\cpge{mpsi-mp2i.nombres-reels-suites.i, pcsi.nombres-reels-suites.g, ptsi.nombres-reels-suites.g, mpsi-mp2i.nombres-reels-suites.d, pcsi.nombres-reels-suites.c, ptsi.nombres-reels-suites.c, mpsi-mp2i.analyse-asymptotique.a, pcsi.analyse-asymptotique.a, ptsi.analyse-asymptotique.a, mpsi-mp2i.procedes-sommatoires-discrets.a, pcsi.series-numeriques.a, ptsi.series-numeriques.a, mpsi-mp2i.procedes-sommatoires-discrets.c, pcsi.series-numeriques.c, ptsi.series-numeriques.c, bcpst1.suites-reelles, bcpst1.limites-continuite.b, bcpst1.derivation.a, mp-mpi.series-numeriques-vectorielles.b, pc.series-numeriques, psi.series-numeriques.a, pt.series-numeriques.a, bcpst2.series-reelles}
\module{Analyse}
\niveau{L1}
\difficulte{}

\contenu{
\question{
On consid{\`e}re la fonction $f$ : $x \mapsto\ln \bigl(\frac{e^x-1}{x}\bigr)$ et la suite d{\'e}finie par
$\begin{cases}u_0\in\R^*\cr u_{n+1}=f(u_n).\cr\end{cases}$
{\'E}tudier la suite $(u_n)$, puis la s{\'e}rie $\sum u_n$.
}
\reponse{
Pour $u_0>0$ on a $u_n\searrow0$ et pour $u_0<0$ on a $u_n\nearrow0$.
$f'(0)=\frac12$ donc $u_{n+1}\sim\frac12u_n$ et la s{\'e}rie $\sum u_n$ converge absolument
(d'Alembert).
}
}
